%%%%%%%%%%%%%%%%%%%%%%%%%%%%%%%%%%%%%%%%%%%%%%%%
% Verbandbuch als Loseblatt: Seite 1 Anschreiben für einen Fensterumschlag
% (DIN 5008 Form A), Seite 2 Dokumentation von Erste-Hilfe-Leistungen.
% Immer beide Seiten als Duplex drucken.
% CC-BY-SA 3.0, FAU FabLab
% Früher Verbandbuch/Verbandbuch-Loseblatt.doc, jetzt LaTeX.
%%%%%%%%%%%%%%%%%%%%%%%%%%%%%%%%%%%%%%%%%%%%%%%%
\documentclass[a4paper,11pt]{scrartcl}
\usepackage{fablab-aushang}
\geometry{left=25mm,right=20mm,top=12mm,bottom=16mm}

\definecolor{ersteHilfe}{HTML}{006A00}
% Strich zum Ausfüllen
\newcommand{\Feld}[1]{\rule{#1}{0.4pt}}
\newcommand{\Zeit}{\Feld{8mm}\,:\,\Feld{8mm} Uhr am \Feld{8mm}\,.\,\Feld{8mm}\,.\,20\Feld{8mm}}

% Seite 1: Kopf, Anschrift und Falzmarken an festen Stellen
\newcommand{\Pos}[3]{\put(\LenToUnit{#1},\LenToUnit{\paperheight-#2}){#3}}
\AddToShipoutPictureFG*{%
  \AtPageLowerLeft{%
    % Kopf: Logo, grauer Balken, Erste-Hilfe-Zeichen
    \Pos{25mm}{24mm}{\colorbox{kopf}{\makebox[155mm][l]{\Logo{11mm}\hfill
      \tikz{\fill[ersteHilfe] (0,0) rectangle (11mm,11mm);
        \fill[white] (4mm,1.8mm) rectangle (7mm,9.2mm) (1.8mm,4mm) rectangle (9.2mm,7mm);}}}}
    % Rücksendeangabe und Anschrift im Fenster
    \Pos{25mm}{33mm}{\scriptsize FAU FabLab · c/o GS EEI · Cauerstr.~7 · 91058 Erlangen}
    \Pos{25mm}{52mm}{\parbox[t]{80mm}{Per Hauspost an:\\
      \textbf{Geschäftsstelle EEI, Sekretariat FAU FabLab}\\
      Cauerstraße 7, 1.~OG, Erlangen}}
    % Angaben rechts
    \Pos{125mm}{38mm}{\parbox[t]{65mm}{%
      Datum:\\[4mm] {\large\Feld{10mm}\,.\,\Feld{10mm}\,.\,20\Feld{10mm}}\\[4mm]
      {\footnotesize Blatt nach der Aufbewahrungsfrist von \emph{fünf} Jahren für
      Verbandbücher vernichten (§\,24 Abs.\,6 DGUV Vorschrift~1).}\\[3mm]
      Blatt-Nummer: {\large\Feld{12mm}\,/\,20\Feld{10mm}}\\
      (wird durch das Sekretariat vergeben)}}
    % Falzmarken und Lochmarke
    \Pos{5mm}{87mm}{\rule{5mm}{0.4pt}}
    \Pos{5mm}{148.5mm}{\rule{8mm}{0.4pt}}
    \Pos{5mm}{192mm}{\rule{5mm}{0.4pt}}}}

\begin{document}
% ------------------------------ Seite 1 ------------------------------
\setlength{\parskip}{0.6em}
\vspace*{78mm}
{\large\bfseries Verbandbuch}

\bigskip
Hallo liebe Nutzerin, hallo lieber Nutzer,

es ist nicht schlimm, wenn du dich bei uns verletzt hast. Material zur Ersten Hilfe findest
du im Verbandschrank am Übergang von Hauptraum und Elektrowerkstatt.

Bitte melde dich bei jemandem von uns Betreuerinnen und Betreuern und fülle die umseitige
Dokumentation von Erste-Hilfe-Leistungen aus (auch bekannt als Verbandbuch). Dieses
Formular steckst du in einen Briefumschlag mit Fenster und lässt es in die Hauspost geben.
Unser Sekretariat wird die Dokumentation datenschutzgerecht aufbewahren.

\textbf{Bitte stelle dich im Zweifel immer einem Arzt vor.}

\bigskip
Grüße und gute Besserung,\\
deine FabLab-Betreuer

\bigskip
\textbf{Anlage}\\
Dokumentation von Erste-Hilfe-Leistungen (auf der Rückseite)

\vfill
\fcolorbox{black}{kopf}{\parbox{\dimexpr\linewidth-2\fboxsep-2\fboxrule\relax}{\itshape
  Hinweis an das Sekretariat:\\
  Bitte dieses Blatt mit der Dokumentation von Erste-Hilfe-Leistungen in den Ordner
  „FabLab-Verbandbuch“ abheften, nachdem oben rechts die laufende Nummer eingetragen wurde.
  Danke!}}

\medskip
{\footnotesize Beim Nachdrucken immer beide Seiten als Duplex ausdrucken.}
\newpage

% ------------------------------ Seite 2 ------------------------------
\setlength{\parskip}{0pt}
\setlength{\tabcolsep}{5pt}
\renewcommand{\arraystretch}{1.25}
\newcommand{\Abschnitt}[1]{\rowcolor{black!55}\multicolumn{1}{|c|}{\color{white}\large #1}\\ \hline}
% \Zeile{Höhe}{Inhalt}
\newcommand{\Zeile}[2]{\rule[-#1]{0pt}{#1}\parbox[t]{\linewidth}{#2}\\ \hline}
\begin{tabularx}{\linewidth}{|X|}
  \hline
  \rowcolor{kopf}\multicolumn{1}{|c|}{\rule{0pt}{6mm}\large\bfseries Dokumentation von Erste-Hilfe-Leistungen}\\
  \rowcolor{kopf}\multicolumn{1}{|c|}{\rule[-2mm]{0pt}{2mm}(nach §\,24 Abs.\,6 DGUV Vorschrift~1)}\\ \hline
  \Abschnitt{Angaben zum Hergang des Unfalls bzw. des Gesundheitsschadens}
  \Zeile{16mm}{Name der verletzten bzw. erkrankten Person:}
  \Zeile{9mm}{Datum und Uhrzeit:\\ \hspace*{40mm}\Zeit}
  \Zeile{24mm}{Ort:\\ \hspace*{3mm}\parbox[t]{80mm}{Räume des FAU FabLab\\ U1.238, U1.239, U1.240\\
    Erwin-Rommel-Straße 60\\ 91058 Erlangen}}
  \Zeile{24mm}{Hergang:}
  \Zeile{24mm}{Art und Umfang der Verletzung bzw. Erkrankung:}
  \Zeile{18mm}{Name der Zeugen:\\[11mm] $\square$ keine Zeugen, nur die verletzte bzw. erkrankte Person}
  \Abschnitt{Erste-Hilfe-Leistungen}
  \Zeile{9mm}{Datum und Uhrzeit:\\ \hspace*{40mm}\Zeit}
  \Zeile{24mm}{Art und Weise der Erste-Hilfe-Maßnahmen:\\[17mm]
    $\square$ nur Wundschnellverband \hfill $\square$ Vorstellung bei Arzt veranlasst bzw. wird erfolgen}
  \Zeile{16mm}{Name der Erste Hilfe leistenden Person:\\[9mm] $\square$ Erste Hilfe an sich selbst geleistet}
\end{tabularx}

\end{document}
